\chapter{Graphs: Topological Sort and Directed Graph Problems}
\chapterepigraph{A directed graph adds a constraint plain traversal
never had to think about: order. Topological sort answers exactly one
question -- in what sequence can these tasks be done so that every
prerequisite comes before what depends on it -- and that single
question, once answered, turns out to be the engine behind cycle
detection, safe-state analysis, and even reconstructing an alien
alphabet from sorted words.}

\section{Ordering a DAG}
\label{sec:graph-topo-intro}

A \textbf{topological sort} of a Directed Acyclic Graph (DAG) is a
linear ordering of its nodes such that for every directed edge
$u \to v$, $u$ appears \textbf{before} $v$ in the ordering. Such an
ordering exists if and only if the graph has \textbf{no cycle} -- a
cycle would force some node to both precede and follow itself,
which is impossible in a linear order.

\begin{patternbox}[Recognizing a Topological Sort Problem]
\emph{``Task scheduling with prerequisites,'' ``compilation order,''
``course ordering,'' or any phrase implying a valid sequence must
respect a set of directed before/after constraints} all point directly
at topological sort.
\end{patternbox}

\begin{ideabox}[Two Ways to Compute It]
\begin{itemize}
  \item \textbf{DFS-based (Section~\ref{sec:graph-topo-sort-dfs})}: run
        DFS; the moment a node has no more unvisited neighbors to
        explore (it is fully ``finished''), push it onto a stack. The
        stack, popped from top to bottom, is a valid topological order
        -- a node finishes \emph{after} everything reachable from it,
        so it naturally ends up before its dependents once the stack
        is reversed.
  \item \textbf{Kahn's Algorithm / BFS-based (Section~\ref{sec:graph-topo-sort-kahn}):}
        repeatedly remove nodes with \textbf{in-degree zero} (no
        remaining unprocessed prerequisites), decrementing the
        in-degree of their neighbors as they're removed -- directly
        simulating ``do everything whose prerequisites are already
        done, then repeat.''
\end{itemize}
\end{ideabox}

\begin{patternbox}[What This Chapter Covers]
\begin{itemize}
  \item \textbf{Topological Sort, both algorithms (Sections~\ref{sec:graph-topo-sort-dfs}--\ref{sec:graph-topo-sort-kahn})}.
  \item \textbf{Detect Cycle in a Directed Graph (Section~\ref{sec:graph-cycle-directed})}:
        a directed-graph adaptation of the parent-tracking idea from
        the BFS/DFS chapter, plus a free cycle check baked into Kahn's
        algorithm itself.
  \item \textbf{Course Schedule I and II (Sections~\ref{sec:graph-course-schedule-1}--\ref{sec:graph-course-schedule-2})}:
        direct applications of cycle detection and topological sort to
        a course-prerequisite framing.
  \item \textbf{Eventual Safe States (Section~\ref{sec:graph-eventual-safe-states})}:
        a three-color DFS variant identifying nodes that can never
        reach a cycle.
  \item \textbf{Alien Dictionary (Section~\ref{sec:graph-alien-dictionary})}:
        building a graph from word-comparison constraints, then
        running topological sort to recover a letter ordering.
\end{itemize}
\end{patternbox}

\newpage

\section{Topological Sort using DFS}
\label{sec:graph-topo-sort-dfs}

\problemstatement{Given a Directed Acyclic Graph (DAG) with $n$ nodes,
return a valid topological ordering -- a linear arrangement where
every edge $u \to v$ has $u$ appearing before $v$.}

\begin{patternbox}
\emph{``Order nodes so every prerequisite comes first''} is solved by
DFS with one addition: push a node onto a \textbf{stack} the moment it
\textbf{finishes} (every neighbor it can reach has already been fully
explored). Popping the stack from top to bottom yields a valid
topological order.
\end{patternbox}

\subsection*{Why finish-time ordering works}

When DFS finishes exploring node $u$ entirely (including everything
reachable from it) before pushing $u$ onto the stack, every node $v$
reachable from $u$ has \emph{already} been pushed (it must finish
before $u$ does, since $u$'s own finish waits on its recursive calls
returning). So $v$ ends up \textbf{lower} in the stack than $u$ -- and
popping top-to-bottom visits $u$ before $v$, exactly the required
order for edge $u \to v$.

\subsection*{Algorithm}

\begin{enumerate}
  \item Maintain a $\textit{visited}$ array and an empty stack.
  \item DFS from every unvisited node: recurse into all unvisited
        neighbors first, and only \textbf{after} that recursion
        returns (the node is now fully finished), push the current
        node onto the stack.
  \item Pop the entire stack to produce the topological order.
\end{enumerate}

\subsection*{C++ implementation}

\begin{lstlisting}
#include <bits/stdc++.h>
using namespace std;

void dfs(int node, vector<vector<int>>& adj, vector<bool>& visited, stack<int>& finishOrder) {
    visited[node] = true;
    for (int neighbor : adj[node]) {
        if (!visited[neighbor]) {
            dfs(neighbor, adj, visited, finishOrder);
        }
    }
    finishOrder.push(node); // push only after all neighbors are done
}

vector<int> topoSortDFS(int n, vector<vector<int>>& adj) {
    vector<bool> visited(n, false);
    stack<int> finishOrder;

    for (int i = 0; i < n; i++) {
        if (!visited[i]) {
            dfs(i, adj, visited, finishOrder);
        }
    }

    vector<int> result;
    while (!finishOrder.empty()) {
        result.push_back(finishOrder.top());
        finishOrder.pop();
    }
    return result;
}

int main() {
    int n = 6;
    vector<vector<int>> adj(n);
    adj[5] = {2, 0}; adj[4] = {0, 1}; adj[2] = {3}; adj[3] = {1};

    for (int x : topoSortDFS(n, adj)) cout << x << " ";
    cout << endl; // 5 4 2 3 1 0  (one valid order)
    return 0;
}
\end{lstlisting}

\subsection*{Dry run}

\dryrun{Take edges $5{\to}2$, $5{\to}0$, $4{\to}0$, $4{\to}1$,
$2{\to}3$, $3{\to}1$.}

DFS from $0$: no outgoing edges, finishes immediately, push $0$.
Stack: $[0]$. DFS from $1$: no outgoing edges, push $1$. Stack:
$[0,1]$. DFS from $2$: recurse to $3$; $3$ recurses to $1$ (already
visited, skip); $3$ finishes, push $3$. Stack: $[0,1,3]$. Back at $2$:
finishes, push $2$. Stack: $[0,1,3,2]$. DFS from $4$: neighbors $0,1$
both visited; finishes, push $4$. Stack: $[0,1,3,2,4]$. DFS from $5$:
neighbors $2,0$ both visited; finishes, push $5$. Stack:
$[0,1,3,2,4,5]$. Popping top to bottom: $5,4,2,3,1,0$ -- every edge
($5{\to}2$, $5{\to}0$, $4{\to}0$, $4{\to}1$, $2{\to}3$, $3{\to}1$)
indeed has its source appearing before its destination.

\begin{complexitybox}
\textbf{Time Complexity.} $O(V+E)$ -- standard DFS traversal cost.
\textbf{Space Complexity.} $O(V)$ for \textit{visited}, the stack, and
the recursion stack.
\end{complexitybox}

\begin{takeawaybox}
``Push on finish, read the stack top-down'' is a clean, purely local
rule -- no node needs any global information about the rest of the
graph to know when to push itself, only that its own DFS call has
returned. This locality is what makes the DFS approach to topological
sort so simple to implement correctly.
\end{takeawaybox}

\section{Topological Sort using Kahn's Algorithm (BFS)}
\label{sec:graph-topo-sort-kahn}

\problemstatement{Same as Section~\ref{sec:graph-topo-sort-dfs}: return
a valid topological ordering of a DAG -- this time using an
in-degree-based BFS instead of DFS finish times.}

\begin{patternbox}
\emph{``Repeatedly do whatever has no remaining prerequisites''} is
solved directly: track each node's \textbf{in-degree} (number of
unprocessed incoming edges); start with every in-degree-$0$ node in a
queue; each time a node is processed, decrement its neighbors'
in-degrees, and any neighbor that drops to $0$ joins the queue.
\end{patternbox}

\subsection*{Why this directly produces a valid order}

A node can only be safely placed in the output once \textbf{every}
prerequisite pointing into it has already been placed -- which is
exactly what in-degree $0$ means at the moment a node is dequeued (all
of its incoming edges have already been ``consumed'' by their source
nodes being processed earlier). This makes Kahn's algorithm a direct
simulation of the real-world process the topological sort is modeling:
do everything that's ready, which unlocks more things to become ready,
repeat.

\subsection*{Algorithm}

\begin{enumerate}
  \item Compute the in-degree of every node by scanning all edges.
  \item Push every node with in-degree $0$ into a queue.
  \item While the queue is non-empty: pop a node, append it to the
        result; for each of its neighbors, decrement their in-degree,
        and push any neighbor whose in-degree reaches $0$.
\end{enumerate}

\subsection*{C++ implementation}

\begin{lstlisting}
#include <bits/stdc++.h>
using namespace std;

vector<int> topoSortKahn(int n, vector<vector<int>>& adj) {
    vector<int> inDegree(n, 0);
    for (int u = 0; u < n; u++) {
        for (int v : adj[u]) inDegree[v]++;
    }

    queue<int> q;
    for (int i = 0; i < n; i++) {
        if (inDegree[i] == 0) q.push(i);
    }

    vector<int> result;
    while (!q.empty()) {
        int node = q.front();
        q.pop();
        result.push_back(node);

        for (int neighbor : adj[node]) {
            inDegree[neighbor]--;
            if (inDegree[neighbor] == 0) q.push(neighbor);
        }
    }
    return result;
}

int main() {
    int n = 6;
    vector<vector<int>> adj(n);
    adj[5] = {2, 0}; adj[4] = {0, 1}; adj[2] = {3}; adj[3] = {1};

    for (int x : topoSortKahn(n, adj)) cout << x << " ";
    cout << endl; // 4 5 2 0 3 1  (one valid order)
    return 0;
}
\end{lstlisting}

\subsection*{Dry run}

\dryrun{Same graph as Section~\ref{sec:graph-topo-sort-dfs}: edges
$5{\to}2,5{\to}0,4{\to}0,4{\to}1,2{\to}3,3{\to}1$.}

In-degrees: $0$ has in-degree $2$ (from $5,4$); $1$ has in-degree $2$
(from $4,3$); $2$ has in-degree $1$ (from $5$); $3$ has in-degree $1$
(from $2$); $4,5$ have in-degree $0$. Initial queue: $[4,5]$. Pop $4$:
result $[4]$; decrement $0\to1$, $1\to1$ (neither reaches $0$ yet).
Pop $5$: result $[4,5]$; decrement $2\to0$ (push $2$), $0\to0$ (push
$0$). Queue: $[2,0]$. Pop $2$: result $[4,5,2]$; decrement $3\to0$
(push $3$). Pop $0$: result $[4,5,2,0]$; no outgoing edges. Pop $3$:
result $[4,5,2,0,3]$; decrement $1\to0$ (push $1$). Pop $1$: result
$[4,5,2,0,3,1]$. Final order: $4,5,2,0,3,1$ -- a different but equally
valid ordering from the DFS approach's $5,4,2,3,1,0$ (topological
order is rarely unique).

\begin{complexitybox}
\textbf{Time Complexity.} $O(V+E)$ -- computing in-degrees scans every
edge once; the BFS itself also visits every node and edge once.
\textbf{Space Complexity.} $O(V)$ for the in-degree array and the
queue.
\end{complexitybox}

\begin{takeawaybox}
Kahn's algorithm has one advantage DFS's approach doesn't offer for
free: if the queue empties before every node has been processed, the
remaining unprocessed nodes are necessarily stuck in a cycle (nothing
in a cycle can ever reach in-degree $0$, since each node in the cycle
is always waiting on another node in the same cycle). This makes Kahn's
algorithm double as a \textbf{cycle detector} in directed graphs, used
directly in the next section.
\end{takeawaybox}

\section{Detect Cycle in a Directed Graph}
\label{sec:graph-cycle-directed}

\problemstatement{Given a directed graph, determine whether it
contains a cycle.}

\begin{patternbox}
Unlike the undirected case (the Graph BFS and DFS Problems chapter's Detect Cycle in an Undirected Graph section),
a single ``parent'' check is not enough here: in a \emph{directed}
graph, reaching an already-visited node is perfectly normal (two
different paths can legitimately converge on the same node without any
cycle). What signals a true cycle is reaching a node that is still
\textbf{on the current DFS path} -- tracked with a second array,
separate from \textit{visited}, marking only nodes currently ``in
progress.''
\end{patternbox}

\subsection*{Why \textit{visited} alone is insufficient here}

Consider $A \to B$, $A \to C$, $B \to D$, $C \to D$: node $D$ is
visited twice (once via $B$, once via $C$), but there is no cycle --
both are legitimate converging paths. The distinguishing feature of an
actual cycle is a \textbf{back edge}: an edge pointing to a node that
is an \textbf{ancestor} of the current node in the DFS recursion tree
(i.e., still actively on the call stack), not merely visited at some
earlier, already-concluded point in the traversal.

\subsection*{Approach 1: DFS with a Path-Tracking Array}

\begin{lstlisting}
#include <bits/stdc++.h>
using namespace std;

bool dfsHasCycle(int node, vector<vector<int>>& adj, vector<bool>& visited,
                  vector<bool>& pathVisited) {
    visited[node] = true;
    pathVisited[node] = true;

    for (int neighbor : adj[node]) {
        if (!visited[neighbor]) {
            if (dfsHasCycle(neighbor, adj, visited, pathVisited)) return true;
        } else if (pathVisited[neighbor]) {
            return true; // back edge to a node still on the current path
        }
    }
    pathVisited[node] = false; // backtrack: leaving this node's path
    return false;
}

bool hasCycleDFS(int n, vector<vector<int>>& adj) {
    vector<bool> visited(n, false), pathVisited(n, false);
    for (int i = 0; i < n; i++) {
        if (!visited[i] && dfsHasCycle(i, adj, visited, pathVisited)) {
            return true;
        }
    }
    return false;
}

int main() {
    int n = 3;
    vector<vector<int>> adj(n);
    adj[0] = {1}; adj[1] = {2}; adj[2] = {0};

    cout << (hasCycleDFS(n, adj) ? "true" : "false") << endl; // true
    return 0;
}
\end{lstlisting}

\begin{complexitybox}[Approach 1 Complexity]
\textbf{Time.} $O(V+E)$.
\textbf{Space.} $O(V)$ for \textit{visited}, \textit{pathVisited}, and
the recursion stack.
\end{complexitybox}

\subsection*{Approach 2: Kahn's Algorithm (BFS)}

As noted at the close of Section~\ref{sec:graph-topo-sort-kahn}, if
Kahn's algorithm processes fewer than $n$ nodes before its queue
empties, the unprocessed nodes are necessarily trapped in a cycle (a
cycle's nodes can never reach in-degree $0$, since each waits on
another node within the same cycle).

\begin{lstlisting}
#include <bits/stdc++.h>
using namespace std;

bool hasCycleKahn(int n, vector<vector<int>>& adj) {
    vector<int> inDegree(n, 0);
    for (int u = 0; u < n; u++) {
        for (int v : adj[u]) inDegree[v]++;
    }

    queue<int> q;
    for (int i = 0; i < n; i++) {
        if (inDegree[i] == 0) q.push(i);
    }

    int processed = 0;
    while (!q.empty()) {
        int node = q.front();
        q.pop();
        processed++;

        for (int neighbor : adj[node]) {
            inDegree[neighbor]--;
            if (inDegree[neighbor] == 0) q.push(neighbor);
        }
    }
    return processed != n; // fewer than n processed => cycle exists
}

int main() {
    int n = 3;
    vector<vector<int>> adj(n);
    adj[0] = {1}; adj[1] = {2}; adj[2] = {0};

    cout << (hasCycleKahn(n, adj) ? "true" : "false") << endl; // true
    return 0;
}
\end{lstlisting}

\subsection*{Dry run}

\dryrun{Take the $3$-cycle $0{\to}1{\to}2{\to}0$.}

\textbf{DFS}: $\textit{dfs}(0)$: visited$[0]=$true, pathVisited$[0]=$true.
Neighbor $1$: unvisited, recurse. $\textit{dfs}(1)$: visited,
pathVisited both true for $1$. Neighbor $2$: unvisited, recurse.
$\textit{dfs}(2)$: visited, pathVisited both true for $2$. Neighbor
$0$: \textbf{already visited, and pathVisited$[0]=$true} (still on the
current path) -- cycle detected, return \texttt{true}.

\textbf{Kahn's}: in-degrees all equal $1$ (each node has exactly one
incoming edge from the cycle). No node starts with in-degree $0$, so
the queue is empty from the start; $\textit{processed}=0 \neq 3=n$ --
cycle detected immediately, without even examining a single edge.

\begin{complexitybox}[Approach 2 Complexity]
\textbf{Time.} $O(V+E)$.
\textbf{Space.} $O(V)$ for in-degrees and the queue.
\end{complexitybox}

\begin{takeawaybox}
Directed-graph cycle detection needs to distinguish ``visited at some
point'' from ``visited and still active on this path'' -- the
\textit{pathVisited} array (cleared on backtrack) is what undirected
cycle detection's simpler parent check generalizes into once edges have
direction. Kahn's algorithm sidesteps this distinction entirely by
checking a global count instead, which is part of why it is often
preferred for this specific check.
\end{takeawaybox}

\section{Course Schedule I}
\label{sec:graph-course-schedule-1}

\problemstatement{Given $\textit{numCourses}$ and a list of
prerequisite pairs $[a,b]$ (meaning course $a$ requires course $b$ to
be completed first), determine whether it is possible to finish all
courses.}

\begin{patternbox}
\emph{``Can all tasks be completed given prerequisite constraints''} is
a direct restatement of Section~\ref{sec:graph-cycle-directed}: it is
possible to finish all courses if and only if the prerequisite graph
(edge $b \to a$ for every pair $[a,b]$, since $b$ must come first)
contains \textbf{no cycle} -- a cycle of prerequisites would mean some
course indirectly requires itself.
\end{patternbox}

\subsection*{Reduction}

Build a directed graph with $\textit{numCourses}$ nodes, adding an edge
$b \to a$ for every prerequisite pair $[a,b]$ (course $b$ must be
completed before $a$, so the edge points from prerequisite to
dependent). The answer is simply
$\neg\,\textit{hasCycle}(\textit{numCourses}, \textit{adj})$, using
either cycle-detection approach from Section~\ref{sec:graph-cycle-directed}
unchanged.

\subsection*{C++ implementation}

\begin{lstlisting}
#include <bits/stdc++.h>
using namespace std;

bool dfsHasCycle(int node, vector<vector<int>>& adj, vector<bool>& visited,
                  vector<bool>& pathVisited) {
    visited[node] = true;
    pathVisited[node] = true;

    for (int neighbor : adj[node]) {
        if (!visited[neighbor]) {
            if (dfsHasCycle(neighbor, adj, visited, pathVisited)) return true;
        } else if (pathVisited[neighbor]) {
            return true;
        }
    }
    pathVisited[node] = false;
    return false;
}

bool canFinish(int numCourses, vector<vector<int>>& prerequisites) {
    vector<vector<int>> adj(numCourses);
    for (auto& p : prerequisites) {
        int a = p[0], b = p[1];
        adj[b].push_back(a); // b must come before a
    }

    vector<bool> visited(numCourses, false), pathVisited(numCourses, false);
    for (int i = 0; i < numCourses; i++) {
        if (!visited[i] && dfsHasCycle(i, adj, visited, pathVisited)) {
            return false;
        }
    }
    return true;
}

int main() {
    vector<vector<int>> prerequisites = {{1, 0}, {0, 1}};
    cout << (canFinish(2, prerequisites) ? "true" : "false") << endl; // false
    return 0;
}
\end{lstlisting}

\subsection*{Dry run}

\dryrun{Take $\textit{numCourses}=2$,
$\textit{prerequisites}=[[1,0],[0,1]]$.}

Edge from pair $[1,0]$: $0 \to 1$ (course $0$ before course $1$).
Edge from pair $[0,1]$: $1 \to 0$ (course $1$ before course $0$).
This is a $2$-cycle: $0 \to 1 \to 0$. Cycle detection immediately
finds the back edge, so $\textit{canFinish}$ returns \texttt{false} --
correctly, since each course requires the other to be completed first,
an impossible circular dependency.

\begin{complexitybox}
\textbf{Time Complexity.} $O(V+E)$, inherited unchanged from cycle
detection.
\textbf{Space Complexity.} $O(V+E)$ for the adjacency list, plus
$O(V)$ for the cycle-detection arrays.
\end{complexitybox}

\begin{takeawaybox}
Recognizing that ``can all tasks be completed'' is a cycle-detection
question in disguise is the entire difficulty of this problem -- once
that reduction is seen, no new algorithm is needed at all.
\end{takeawaybox}

\section{Course Schedule II}
\label{sec:graph-course-schedule-2}

\problemstatement{Given $\textit{numCourses}$ and prerequisite pairs
$[a,b]$ as in Section~\ref{sec:graph-course-schedule-1}, return
\textbf{one valid order} to take all courses, or an empty array if no
valid order exists.}

\begin{patternbox}
\emph{``Return a valid order respecting prerequisites''} is a direct
restatement of topological sort
(Sections~\ref{sec:graph-topo-sort-dfs}--\ref{sec:graph-topo-sort-kahn}):
build the same prerequisite graph as Course Schedule I, then run either
topological sort algorithm. Kahn's algorithm is the natural choice here
since it detects the no-valid-order case (a cycle) \textbf{for free},
via the same processed-count check from Section~\ref{sec:graph-cycle-directed}.
\end{patternbox}

\subsection*{Reduction}

Build the graph identically to Course Schedule I ($b \to a$ for every
pair $[a,b]$), run Kahn's algorithm, and return the resulting order --
but only if every node was processed; otherwise return an empty array
(a cycle means no valid order exists).

\subsection*{C++ implementation}

\begin{lstlisting}
#include <bits/stdc++.h>
using namespace std;

vector<int> findOrder(int numCourses, vector<vector<int>>& prerequisites) {
    vector<vector<int>> adj(numCourses);
    vector<int> inDegree(numCourses, 0);

    for (auto& p : prerequisites) {
        int a = p[0], b = p[1];
        adj[b].push_back(a);
        inDegree[a]++;
    }

    queue<int> q;
    for (int i = 0; i < numCourses; i++) {
        if (inDegree[i] == 0) q.push(i);
    }

    vector<int> order;
    while (!q.empty()) {
        int node = q.front();
        q.pop();
        order.push_back(node);

        for (int neighbor : adj[node]) {
            inDegree[neighbor]--;
            if (inDegree[neighbor] == 0) q.push(neighbor);
        }
    }

    if ((int)order.size() != numCourses) return {}; // cycle: no valid order
    return order;
}

int main() {
    vector<vector<int>> prerequisites = {{1, 0}, {2, 0}, {3, 1}, {3, 2}};
    vector<int> order = findOrder(4, prerequisites);
    for (int x : order) cout << x << " ";
    cout << endl; // 0 1 2 3  (one valid order)
    return 0;
}
\end{lstlisting}

\subsection*{Dry run}

\dryrun{Take $\textit{numCourses}=4$,
$\textit{prerequisites}=[[1,0],[2,0],[3,1],[3,2]]$ (course $3$ needs
both $1$ and $2$; both $1$ and $2$ need $0$).}

Edges: $0\to1,0\to2,1\to3,2\to3$. In-degrees: $0$:$0$, $1$:$1$,
$2$:$1$, $3$:$2$. Initial queue: $[0]$. Pop $0$: order $[0]$; decrement
$1\to0$ (push), $2\to0$ (push). Queue: $[1,2]$. Pop $1$: order
$[0,1]$; decrement $3\to1$ (not $0$ yet). Pop $2$: order $[0,1,2]$;
decrement $3\to0$ (push). Pop $3$: order $[0,1,2,3]$. Size $4 =
\textit{numCourses}$: valid order found.

\begin{complexitybox}
\textbf{Time Complexity.} $O(V+E)$, inherited unchanged from Kahn's
algorithm.
\textbf{Space Complexity.} $O(V+E)$ for the adjacency list and
in-degree array.
\end{complexitybox}

\begin{takeawaybox}
Course Schedule I and II together demonstrate the cleanest possible
illustration of this chapter's two core tools applied side by side:
the exact same prerequisite graph answers ``is it possible'' via cycle
detection and ``what is an order'' via topological sort, with no
additional derivation needed for either once the graph is built.
\end{takeawaybox}

\section{Eventual Safe States}
\label{sec:graph-eventual-safe-states}

\problemstatement{Given a directed graph, a node is \textbf{safe} if
every possible path starting from it eventually leads to a
\textbf{terminal node} (a node with no outgoing edges) -- meaning no
path from it can ever enter a cycle. Return all safe nodes, sorted.}

\begin{patternbox}
\emph{``Never gets stuck in a cycle, on any path''} extends
Section~\ref{sec:graph-cycle-directed}'s \textit{pathVisited} idea into
a full \textbf{three-state} DFS: a node is currently
\textbf{being explored} (on the active path -- analogous to
\textit{pathVisited}), \textbf{fully confirmed safe}, or
\textbf{not yet visited}. Reaching a node still being explored signals
a cycle, exactly as before -- the new piece is caching the confirmed
\texttt{safe}/\texttt{unsafe} verdict for every node once decided, so
each node's full reachability is resolved only once.
\end{patternbox}

\subsection*{Why a single boolean \textit{visited} array is not enough}

A node might be reached by DFS, fully resolved as safe or unsafe, and
then reached \emph{again} via a different path later in the traversal
-- a plain \textit{visited} array would either skip re-deriving its
answer (fine, if it remembers the answer) or, worse, treat it as
``currently on the path'' if visited and not-yet-finished are
conflated into one state. Three explicit states --
$0=\textit{unvisited}$, $1=\textit{visiting}$, $2=\textit{safe}$ --
avoid this ambiguity entirely.

\subsection*{Algorithm}

\begin{enumerate}
  \item For each node, maintain a state: $0$ (unvisited), $1$
        (currently being explored -- on the active DFS path), or $2$
        (confirmed safe).
  \item $\textit{dfs}(\textit{node})$: if state is $2$, return
        \texttt{true} (already confirmed safe). If state is $1$,
        return \texttt{false} (currently on the path -- this is a back
        edge, a cycle). Otherwise, mark state $1$ (entering); recurse
        into every neighbor -- if \emph{any} neighbor's DFS returns
        \texttt{false}, this node is unsafe, return \texttt{false}
        immediately. If every neighbor returns \texttt{true}, mark
        state $2$ (confirmed safe) and return \texttt{true}.
  \item Collect every node whose DFS returns \texttt{true}.
\end{enumerate}

\subsection*{C++ implementation}

\begin{lstlisting}
#include <bits/stdc++.h>
using namespace std;

bool dfs(int node, vector<vector<int>>& adj, vector<int>& state) {
    if (state[node] == 2) return true;
    if (state[node] == 1) return false;

    state[node] = 1; // currently being explored
    for (int neighbor : adj[node]) {
        if (!dfs(neighbor, adj, state)) return false;
    }
    state[node] = 2; // confirmed safe
    return true;
}

vector<int> eventualSafeNodes(int n, vector<vector<int>>& adj) {
    vector<int> state(n, 0);
    vector<int> result;

    for (int i = 0; i < n; i++) {
        if (dfs(i, adj, state)) {
            result.push_back(i);
        }
    }
    return result; // already sorted, since nodes are processed in order
}

int main() {
    int n = 7;
    vector<vector<int>> adj(n);
    adj[0] = {1, 2}; adj[1] = {2, 3}; adj[2] = {5}; adj[3] = {0};
    adj[4] = {5}; adj[5] = {}; adj[6] = {};

    for (int x : eventualSafeNodes(n, adj)) cout << x << " ";
    cout << endl; // 2 4 5 6
    return 0;
}
\end{lstlisting}

\subsection*{Dry run}

\dryrun{Take the graph above: $0\to1,0\to2,1\to2,1\to3,2\to5,3\to0,
4\to5$ ($5,6$ terminal).}

$\textit{dfs}(0)$: state$[0]=1$. Neighbor $1$: $\textit{dfs}(1)$:
state$[1]=1$. Neighbor $2$: $\textit{dfs}(2)$: state$[2]=1$. Neighbor
$5$: $\textit{dfs}(5)$: no neighbors, state$[5]=2$, return
\texttt{true}. Back at $2$: state$[2]=2$, return \texttt{true}. Back
at $1$: neighbor $3$: $\textit{dfs}(3)$: state$[3]=1$. Neighbor $0$:
state$[0]=1$ -- \textbf{still being explored} -- return
\texttt{false}. So $\textit{dfs}(3)$ returns \texttt{false}, making
$\textit{dfs}(1)$ return \texttt{false} (it has an unsafe neighbor),
which in turn makes $\textit{dfs}(0)$ return \texttt{false}. Continuing
with $i=2$: already state $2$, returns \texttt{true}, added. $i=4$:
$\textit{dfs}(4)$: neighbor $5$, already safe (state $2$); $4$ becomes
safe. $i=5,6$: already safe / trivially safe (no neighbors). Final
safe nodes: $2,4,5,6$ -- matching the expected output, while $0,1,3$
are correctly excluded since they all lie on or lead into the cycle
$0\to1\to3\to0$.

\begin{complexitybox}
\textbf{Time Complexity.} $O(V+E)$ -- each node's DFS result is
computed and cached exactly once (the state-$2$ check short-circuits
any repeat work).
\textbf{Space Complexity.} $O(V)$ for the state array and the
recursion stack.
\end{complexitybox}

\begin{takeawaybox}
The white/gray/black three-coloring scheme used here is the general
form of cycle-aware DFS: ``gray'' (in-progress) plays the role
\textit{pathVisited} played in plain cycle detection, while ``black''
(confirmed done) adds memoization on top, caching each node's final
verdict so a graph with many converging paths is still only ever
analyzed once per node.
\end{takeawaybox}

\section{Alien Dictionary}
\label{sec:graph-alien-dictionary}

\problemstatement{Given a list of words sorted lexicographically
according to the rules of an unknown alien language with $k$
distinct lowercase letters, determine a valid ordering of the
alphabet. Return an empty string if no valid ordering exists.}

\begin{patternbox}
This closing problem of the chapter is a two-step composition:
\textbf{build} a directed graph of letter-ordering constraints by
comparing each pair of \textbf{adjacent} words in the list, then
\textbf{run topological sort} (Sections~\ref{sec:graph-topo-sort-dfs}--\ref{sec:graph-topo-sort-kahn})
on that graph to recover the alphabet order. The graph-building step is
the only new idea; the ordering step is unchanged from earlier in this
chapter.
\end{patternbox}

\subsection*{Extracting constraints from adjacent word pairs}

For two adjacent words $w_1, w_2$ in the sorted list, scan character by
character until the \textbf{first} position where they differ. If
$w_1[i] \neq w_2[i]$, that single comparison reveals exactly one
constraint: $w_1[i]$ comes before $w_2[i]$ in the alien alphabet
(edge $w_1[i] \to w_2[i]$) -- and nothing more can be concluded from
this pair, since the relative order of any later characters is
unconstrained by this comparison. If no differing position is found
within the shorter word's length, the ordering is \textbf{invalid}
exactly when $w_1$ is \emph{longer} than $w_2$ (a longer word cannot
correctly precede its own prefix in sorted order).

\subsection*{Algorithm}

\begin{enumerate}
  \item Initialize an empty adjacency structure over the (at most
        $26$) distinct letters appearing in the words.
  \item For every adjacent pair of words, find the first differing
        character and add the corresponding edge; if no differing
        character exists and the first word is longer, return
        \texttt{""} immediately (invalid).
  \item Run topological sort (Kahn's algorithm is convenient here,
        since it naturally detects a cycle -- contradictory
        constraints -- by an incomplete processed count) over the
        graph of letters.
  \item If a cycle is detected, return \texttt{""}; otherwise return
        the topological order as a string.
\end{enumerate}

\subsection*{C++ implementation}

\begin{lstlisting}
#include <bits/stdc++.h>
using namespace std;

string alienOrder(vector<string>& words) {
    unordered_set<char> letters;
    for (string& w : words) for (char c : w) letters.insert(c);

    unordered_map<char, unordered_set<char>> adj;
    unordered_map<char, int> inDegree;
    for (char c : letters) inDegree[c] = 0;

    for (int i = 0; i + 1 < (int)words.size(); i++) {
        string& w1 = words[i];
        string& w2 = words[i + 1];
        int minLen = min(w1.size(), w2.size());
        bool foundDiff = false;

        for (int j = 0; j < minLen; j++) {
            if (w1[j] != w2[j]) {
                if (adj[w1[j]].find(w2[j]) == adj[w1[j]].end()) {
                    adj[w1[j]].insert(w2[j]);
                    inDegree[w2[j]]++;
                }
                foundDiff = true;
                break;
            }
        }
        if (!foundDiff && w1.size() > w2.size()) return "";
    }

    queue<char> q;
    for (char c : letters) {
        if (inDegree[c] == 0) q.push(c);
    }

    string order;
    while (!q.empty()) {
        char c = q.front();
        q.pop();
        order += c;

        for (char neighbor : adj[c]) {
            inDegree[neighbor]--;
            if (inDegree[neighbor] == 0) q.push(neighbor);
        }
    }

    return (int)order.size() == (int)letters.size() ? order : "";
}

int main() {
    vector<string> words = {"wrt", "wrf", "er", "ett", "rftt"};
    cout << alienOrder(words) << endl; // wertf
    return 0;
}
\end{lstlisting}

\subsection*{Dry run}

\dryrun{Take $\textit{words}=[\texttt{wrt},\texttt{wrf},\texttt{er},
\texttt{ett},\texttt{rftt}]$.}

Pair (\texttt{wrt},\texttt{wrf}): differ at index $2$ ($\texttt{t}$
vs $\texttt{f}$) -- edge $t \to f$. Pair (\texttt{wrf},\texttt{er}):
differ at index $0$ ($\texttt{w}$ vs $\texttt{e}$) -- edge $w \to e$.
Pair (\texttt{er},\texttt{ett}): differ at index $1$ ($\texttt{r}$ vs
$\texttt{t}$) -- edge $r \to t$. Pair (\texttt{ett},\texttt{rftt}):
differ at index $0$ ($\texttt{e}$ vs $\texttt{r}$) -- edge $e \to r$.
Graph: $t\to f$, $w\to e$, $r\to t$, $e\to r$. In-degrees: $w$:$0$,
$e$:$1$, $r$:$1$, $t$:$1$, $f$:$1$. Kahn's: start $[w]$. Pop $w$:
order $\texttt{w}$; $e\to0$, push. Pop $e$: order $\texttt{we}$;
$r\to0$, push. Pop $r$: order $\texttt{wer}$; $t\to0$, push. Pop $t$:
order $\texttt{wert}$; $f\to0$, push. Pop $f$: order
$\texttt{wertf}$. All $5$ letters processed: valid order
$\texttt{"wertf"}$.

\begin{complexitybox}
\textbf{Time Complexity.} $O(\sum |w_i| + V + E)$ -- comparing each
adjacent word pair costs time proportional to the shorter word's
length, and the topological sort afterward costs $O(V+E)$ over the
(at most $26$) letters and their constraints.
\textbf{Space Complexity.} $O(V+E)$ for the letter graph, bounded by
$O(26^2)$ in the worst case.
\end{complexitybox}

\begin{takeawaybox}
This chapter closes by demonstrating topological sort's full
generality: the ``nodes'' and ``edges'' don't need to be given
explicitly at all -- here they are \emph{derived} from a completely
different kind of input (a sorted word list) via a simple
pairwise-comparison rule, and once that translation step is done, every
algorithm from earlier in this chapter (cycle detection bundled into
Kahn's algorithm, the topological order itself) applies completely
unchanged.
\end{takeawaybox}

